\section{Conclusioni}
\label{sec:conclusions}

Gli esperimenti mostrano che il transfer learning con ResNet18 e la conservazione del
dettaglio spaziale sono più utili delle loss specializzate provate. Mean Dice e PQ
misurano aspetti diversi: una previsione più selettiva può riconoscere meglio le istanze,
ma ridurre la sovrapposizione della maschera complessiva.

L'esperimento \finalexperiment{} è il modello consegnato: una U-Net ResNet18 addestrata
su crop casuali a $1536^2$ da frame nativi a $2048^2$. Ottiene
\textbf{\finaldice\ di Mean Dice} e \textbf{\finalpq\ di PQ}; il picco GPU è
\finaltrainmemory{} in training e \finaltestmemory{} in inferenza.

La configurazione 15m raggiunge 0.6650 di Mean Dice e 0.4153 di PQ entro i limiti di
memoria, mentre i modelli più pesanti aumentano la PQ ma non il Mean Dice e richiedono
molte più risorse.

Il limite principale è il possibile ottimismo dovuto al riuso della validation per il
confronto e la selezione dell'output; in alcune prove il cambio del seed ha comunque
prodotto risultati pressoché invariati. Un possibile sviluppo è migliorare il recall dei
filamenti più piccoli senza aumentare i falsi positivi.
